\section{Year in Review Continued}

\begin{quotex}
The direction in which we may find awakening and liberation, the direction of liberation, is vertical and has nothing to do with the course of history.
\flright{\textsc{Julius Evola}, \emph{The Doctrine of Awakening}}
\end{quotex}


Rather than respond to comments, I'll address several issues here. As I've indicated several times, our goal at Gornahoor has been to prepare an intellectual elite. Over time, the audience has changed quite a bit to the point where there is now more of a common mind or at least a general agreement of principles. In post 1001, I will reveal everything.

The end of my involvement in Gornahoor is not necessarily the end of Gornahoor. As more writers become involved, it will develop a different identity. It could take some time, perhaps generations, for the elite to become effective; after all, the time of preparation began some 90 years ago. Some readers think I have the alchemical skill to translate spiritual ideas to the written text and they have benefited from it. That was the intent: to clear away false ideas to get to the core of Tradition. Once certain principles are understood, the ensuing choice and commitment becomes clear and necessary.

Ideally, there should be 100 new blogs bringing Tradition to the various walks of life, professions, arts, and crafts. So, I am not predicting an end, but rather a beginning.

\paragraph{Buddhism}


Readers must take care in what they read and not jump to conclusions that do not follow. I try to be precise. Furthermore, texts must be understood in the light of other texts. Reading in this manner is difficult; however, writing in such a way is even more difficult. I hope I have not failed completely. For example, I wrote:

\begin{quotex}
I myself have been initiated into Mahayana Buddhism, since that option was available to me.
\end{quotex}


Given what we have written about initiation and conversion, particularly in the case of Guenon, who can conclude from this that I am a Buddhist? It is obvious that Guenon and Evola have put great stock into the idea of ``initiation''; as the correspondence show, Evola pestered Guenon with that question even 25 years after their first exchange of letters. So I, too, took it quite seriously. I also took seriously Guenon's claim in the King of the World that the `lost word' would be found among the sages of Tibet.

Hence, I studied everything Tibetan, including occultists like Blavatsky, Nicholas Roerich, even Alice Bailey who claimed a Tibetan connection. I took the sangha vows and a few initiations from various Tibetan Lamas. Now that settled the problem of initiation for me so I could continue my studies of Tradition knowing I was an initiate. Was that necessary? Readers can decide for themselves. As for myself, I have been thinking lately ``no''. On the other hand, maybe personal instruction in forms of meditation and hearing doctrine from initiates really has made an effective difference in understanding; that can't be ruled out. In retrospect, it did lead to a change of direction in my life. However, I don't recommend this for all, especially those who have returned to the Western Tradition since they would have to confess that apostasy.

In any case, that did not make me a Buddhist; just an initiate in a valid Tradition. I left when I realized that my purposes had been served. In any case, Buddhism in America is in a sorry state. For example, I received a booklet about my dharma vows, in which there is a general prohibition against oral, anal, and masturbatory sex. Since a large number of American Buddhists are ex-Catholics who have told me they have moved beyond that, these prohibitions are mostly ignored. Even more curiously, a minor transgression against the vows can be repaired with a personal prayer, but a more serious one requires the retaking of the vows with a lama. This is identical to the Catholic distinction between venial and mortal sins.

What I liked about Tibet was that it was the last of the patriarchal theocratic societies, so it is still a good model. Unfortunately, no one else was interested in that sort of Tradition. Apparently, the Dalai Lama himself is thoroughly Westernized, despite his reactionary past if anything on the Internet about him can be believed. A blessing did come my way, from a Jewish convert to Buddhism who had a PhD in medieval history. She convinced me that the Medieval era was much more than the ``Dark Ages'', which gradually altered my view of that time. More detailed study led me to my current understanding.

As for the question about the Doctrine of Awakening, Buddhism has become a ``devotional faith'' by Aryan westerners. I oppose ``large'' judgments about movements that evolve over the centuries involving hundreds of millions of people. They give the impression of erudition while lacking all content. You can judge the Doctrine of Awakening by how well Evola extracts the various inner states of Buddhist practices. He does that well through his quotes and organization of key ideas. Nevertheless, that is merely academic exercise without making the efforts of actually achieving those states. They do not arise, I can assure you, from reading about them.

\paragraph{Priests and Kings Redux}


To return to this tedious theme once again, despite the many times it has come up, no one ever addresses anything specific in those posts. In the first edition of Revolt against the Modern World, Evola quoted a text from the Aitareya Brahmana and claimed that the King was the Sun and the Priest the Moon. Coomaraswamy, a Sanskrit scholar, pointed out Evola's error in a review. The text really said the exact opposite; the Priest recited the bridal vow that an Indian groom makes to his bride. That is the proper relationship between the Priest, as the Sun, and the King as the Moon. Evola pulled that reference from future editions of Revolt, but never changed his thesis. Draw your own conclusions about hiding contrary evidence. From now on, any comments on this issue will have to address that specific text from the Brahmana.

From a metaphysical point of view, it makes sense. The Priest is the link between the physical and the transcendent, i.e., the vertical. The King acts horizontally, although accepting the advice and the counsel of his priests. The vertical is Purusha, the unchanging, the masculine. Action is Prakriti; even in Tantra, Shakti is the feminine principle of action. To twist this to mean something else is a waste of intellectual capital. Together, the vertical and the horizontal form the sign of the cross, the sign of Triumph.

\paragraph{The Western Tradition}


So this brings us to the self-loathing that marks certain segments of the counter-Tradition in the West. They look everywhere for Tradition except for where it is. For us, that is the pre-schism Church, uniting the East and the West horizontally as well as the connection to the best of pagan Greece and Rome. We also need to mention the contribution of the Celts and the Franks to this synthesis.

This includes an esoteric Tradition that has been poorly preserved and also the Hermetic tradition that is also preserved and developed. Note that this is not the same as so-called ``Traditional'' elements in churches today, since they are usually suspicious of anything called esoteric, perennial, or Hermetic. Nevertheless, the esoteric teachings are not intended to replace the exoteric teachings; rather they are reserved for the few who interest themselves in such things.

In modern times, we have mentioned a particular trend that shows the initiation still exists in the West, as Guenon suspected. The Russian esoterist \textbf{G. O. Mebes} initiated both \textbf{Valentin Tomberg} and \textbf{Mouni Sadhu}\footnote{15 Jan 2021: Correction. Tomberg was not initiated by Mebes, but was in direct contact with some of his students.}. We have mentioned less about the latter. Sadhu was actually a Polish Catholic who lived in \textbf{Ramana Maharshi}'s ashram for a time before settling in Australia. Based on his own experiences and studies, he claims that the highest teachings of the West were not less than those of the East. He published several useful books of exercises on concentration, meditation, theurgy, and control of the mind. Some may find them helpful.

Otherwise, you are all on your own, some doing better than others.

\flrightit{Posted on 2013-01-15 by Cologero}

\begin{center}* * *\end{center}

\begin{footnotesize}
\begin{sffamily}
\texttt{Janus on 2013-01-16 at 01:39 said:}

Where in the west would you say one can find the pre-schism Church? In my journey, I have investigated both Catholicism and Orthodoxy. In Catholic churches today, one seems to be faced with a grim choice of a modernizing Church more concerned with condoms in Africa than transforming sacraments on the one hand, and a traditional Catholicism on the other which so often is filled with reactionary elements (in the sense of being hostile to anything perennial or esoteric, as you said). Orthodoxy seems to me to have far more promise, but as Fr. Rose found, hostility can come from this direction as well.

Personally, I stumbled upon Tradition as one not believing in God (a conclusion reached after the journey mentioned above). Investigating the Traditionalist writings, the idea of God as the ``face'' of the Divine provides some reconciliation here to me. As I once read it described, ``God is a mask. Who wears Him?'' The questions of Traditional metaphysics and personal practice are great ones indeed, even without considering the question of God in the theistic conception of Him. As such, Buddhism has increasingly seemed interesting to me, as it is not a theistic path. Would you say that I should look back to Catholicism or Orthodoxy despite this (as thoughtful advice, not guidance from a guru, of course)? One might have either ``a shaykh or Shaitan'' as ones' guide, but when shaykhs seem impossible to find in the remains of the the western Tradition, one really seems left to ones' own devices.

That said, I was reading earlier today Prince Charles' opening address to Sacred Web's 2006 conference. If the Prince of Wales himself is struggling with these questions, perhaps that provides some comfort that we are not so alone as it may appear, and hope that we may indeed see the ``end to all our exploring''. Thank you very much, Cologero, for the work you did here. I particularly enjoyed the Maurras translations, and would definitely like to see more if you are able. For Maurras and de Maistre alone, French is a language to learn. Best wishes, and I hope to see this blog and your work continue after you finish.

\hfill

\texttt{Avery on 2013-01-16 at 02:09 said:}

Janus,

Inherent in the juxtaposition of tradition and modernity is the recognition that things are at a very decayed state today. But I see the existence of this blog as a sign of hope. Guenon diagnosed the disease before the victim was entirely dead. After him there have not been many doctors with the intellectual ability to start treatment. This is the beginning of that era. It is natural, following the intellectual rigor of Guenon, to see decay all around you. Embrace it and correct it with compassion.

\hfill

\texttt{Brett Stevens on 2013-01-16 at 07:35 said:}

Tradition seems to me to be a truth derived from life itself that can be spoken in many languages or philosophical types. However, that doesn't mean it bends like a young tree in the wind. It is roughly the same thing because the world and cosmos are the same everywhere.

I hope you continue writing in whatever form suits you. From having written online for over 20 years, I recommend taking regular breaks anyway. You can simply deplete yourself of what you wish to communicate, and in doing so, neglect internal communication that helps you develop clarity of your own ideas.

\hfill

\texttt{Cologero on 2013-01-16 at 08:24 said:}

I'm afraid, Janus, that the only thing I can tell you is to not confuse a particular institution with the spiritual tradition itself, which is much larger. Also, practices which are suitable for monks are not always suitable for men in the world.

\hfill

\texttt{Visionsofglory14 on 2013-01-17 at 00:48 said:}

You are correct that actual initiation in the modern west might be helpful but is unnecessary and often detrimental if you get involved with the wrong institution or subgroup of an institution. In this Kali Yuga, all of the once carefully guarded sacred texts are easily available for everyone, so once a seeker has the general idea, initiation is perhaps even frivolous. That isn't to say guidance from a master won't be helpful. There's a ton to be misled by out there.

However, we can expect rot from top to bottom in these organizations, so self-education/direction is necessary. Former traditional bulwarks have completely lost their orientation even to the point of inversion of principle. Only in rare glimpses can the former spirit be witnessed. It only remains in disparate individuals. And you can often find that you have a better chance of communicating principle to any half-way decent person you meet than a member of an apparently traditional institution, so restoration is a red herring. Better to simply aim at self-development while communicating to those who cross your path than hoping for a restoration.

\hfill

\texttt{Visionsofglory14 on 2013-01-17 at 01:07 said:}

Doctrine of Awakening is an exemplary work. I don't know where I'd be without it. Evola and Nietzsche are two poles of the same spirit. Nietzsche alone forms a closed circuit as modern hands have neutered him, but Evola is like an outlet that makes spiritual combat possible again. Guenon is the opposite extreme from Nietzsche. Nietzsche=pure activity, Guenon=pure spirit. Evola is an attempt to mediate the two.

\hfill

\texttt{Cologero on 2013-01-17 at 07:39 said:}

How exactly do you ``mediate'' between ``pure spirit'' and ``pure activity''? You say Evola ``attempted'' to perform that mediation ... did he succeed or fail? If the former, how did he do it?


Guenon alone of the three understood the relationship between Purusha (spirit, masculine) and Prakriti (activity, feminine). Nietzsche seemed to intuit this when he wrote, ``when you go to a woman, bring a whip,'' but he probably never understood its true significance. To wit, activity must be guided by spirit, or, a fortiori, activity must be dominated by spirit.

For Nietzsche, there is no spirit and no pole, rather only power and appearance. That may explain why, judging from my facebook feed, Nietzsche appeals primarily to women and homosexuals.

\hfill

\texttt{Jason-Adam on 2013-01-17 at 08:24 said:}

1. I am still surprised that anyone still can admire a man like Nietzsche on the same level as a true man of wisdom such as Guenon or even Evola.

2. Cologero has done an amazing service in pointing us to Mebes, Tomberg and the others as they may be well the long sought after Christian path to transcendence that was occluded after the Middle Ages.

\hfill

\texttt{doctoroctagon on 2013-01-17 at 08:48 said:}

``our goal at Gornahoor has been to prepare an intellectual elite.''

I was put in mind of this: ``Never doubt that a small group of thoughtful, committed, citizens can change the world. Indeed, it is the only thing that ever has.''

Margaret Mead. Also my blog over at \url{http://www.golgonooza.blogspot.com} is going to have some more

Tradition and the arts oriented posts soon -- it might become one of those 100 blogs that you mention!

\hfill

\texttt{Visionsofglory14 on 2013-01-17 at 10:32 said:}

Evola, Guenon, and Nietzsche are united in that they recognize the decadence of modernity and that it must be destroyed so that life can regenerate. The greatest difference between Guenon and Evola is that Guenon sees the Kali Yuga as a blight, even though he recognizes as a necessary step in the return to the Golden Age. Nietzsche, on the other hand, teaches to enter the nihilistic phase with glee.

Nietzsche gave his bible Thus Spoke Zarathustra the subtitle ``A Book for All and None.'' A book that is meant for everyone will be vastly read, read at many different levels, yet clearly, as Nietzsche expected, to be greatly misunderstood. So Cologero `s attempt to defame Nietzsche as a tool for homosexuals and women is sidestepped by Nietzsche's quite obvious knowledge and intention that he would be used for destructive purposes by devious people.

Both Nietzsche and Guenon, however, do see envision a return to greater times. Unknowingly or not, they share this. Guenon's focus is on the preservation of tools that will bring back the golden age and be used by the future elite. Nietzsche, though, recognized that the negative polarity is necessary to divide mankind and make us hopelessly unequal, to the socialist's chagrin. Thus, with his books, he offers both high and low a new found freedom in power. For the high, Evola and Guenon may offer something: Ride the Tiger. For the low, this cannot mean anything more than the tools for their own destruction.

\hfill

\texttt{Cologero on 2013-01-17 at 12:44 said:}

Dear Delusions of Glory: First of all, I did not attempt to ``defame'' anything. I simply stated a fact. A fact is true or false; if you think it is false, then simply refute it. Others will make up their minds about who is ``deviously'' misusing Nietzsche.

Let me remind you, as well as everyone else, we are interested in discussing principles here, not in idle displays of pseudo-erudition of the type found in Spark Notes. So what exactly are the principles embraced by Guenon and Nietzsche? Obviously, Nietzsche claimed that the Will to Power is the fundamental principle of the world. That is, there is no transcendence, nothing hidden, just multiple manifestations of the will to power. There is nothing of Guenon in that, nothing of Tradition in that. Now to mediate --- something I asked you to explain, although you are not quite up to the task --- between Guenon and Nietzsche can mean nothing other than finding a higher principle that would incorporate and encompass both of them. So what is that mediating higher principle? Take your time answering ... You say Evola tried to do that. What he did was distinguish between the best and the worst of Nietzsche, as he called it, much like a boy how tries to pick the peas and carrots out from his beef stew. The question then becomes, what exactly is left of Nietzsche? Is Evola being devious and destructive by ``correcting'' Nietzsche? The same could be said of Coomaraswmay, which we documented here. Yes, we can rely on Nietzsche for his powerful poetry and heroic attitude, but is AKC also devious? Again, take your time, because we expect specific quotes and examples, not empty hand waving.

As for Nietzsche and glee ... does he truly demonstrate that in his life?

\hfill

\texttt{Mihai on 2013-01-18 at 06:47 said:}

If Nietzsche's end goals are taken seriously (instead of merely adopting some good points and means he presented) than he can only be a catalyst for one thing: the counter-tradition.

\hfill

\texttt{Visionsofglory14 on 2013-01-18 at 08:46 said:}

Explain Nietzsche's, Evola's, and Guenon's end goals, since you speak on such high authority.

\hfill

\texttt{Mihai on 2013-01-18 at 09:49 said:}

Explain the difference between Nietzsche's ``superman'' and Marx's ``end of history''- if there are any

(Hints: They are both relativistic, utopic, and they begin and end in the immanent, meaning that both are completely horizontal, lacking in any transcendence. )

\hfill

\texttt{Cologero on 2013-01-19 at 09:31 said:}

I have no obligation to explain anything to you, visions of glory. Rather, you have the obligation to properly prepare yourself before entering comments. Specifically, you need to refer to the actual content of posted texts and especially you need to respond to direct questions. The only authority comes from truth. If there is no truth, there is no authority. That is the fundamental issue and always the starting point. If you deny the existence of truth, or the ability to know it, then there is no independent objective authority. There is only the personal will to power to impose opinions through force.

\hfill

\texttt{andrewlmt on 2013-01-20 at 02:21 said:}

Reading the Church Fathers is good medicine. They are a way ``in''; I believe worthy guides for the Christian.

\hfill

\end{sffamily}
\end{footnotesize}

